% This is formatted for the HWID 2026 Conference using the Springer LNCS template
\documentclass[runningheads]{llncs}
\usepackage[T1]{fontenc}
\usepackage{graphicx}
\usepackage{url}

\begin{document}

\title{Treating Workers Like Guided Missile Pigeons: Cybernetics, Behaviorism, and Symbolic Systems of Control}
\titlerunning{Treating Workers Like Guided Missile Pigeons}

\author{Watson Hartsoe}
\authorrunning{W. Hartsoe}
\institute{Digital Media, Georgia Institute of Technology \\
\email{whartsoe3@gatech.edu}}

\maketitle

\begin{abstract}
Industry 5.0 promises the harmonisation of human and machine intelligence, shifting the workplace focus from pure productivity to human well-being and augmentation \cite{breque2021industry}. However, this paper argues that the Fifth Industrial Revolution (5IR) does not transcend earlier paradigms of behavioral control; rather, it internalizes them. Tracing the genealogy of behavioral optimization from B.F. Skinner's Project Pigeon and Norbert Wiener's cybernetics to contemporary AI-driven workplace analytics, we conceptualize the modern iteration of ORCON (Organic Control). Originally a mid-century military architecture that integrated living organisms into cybernetic loops as servomechanisms, the contemporary ORCON regulates human behavior by co-opting Clifford Geertz's concept of ``thick description.'' By analyzing Operator Digital Twins and Reinforcement Learning from Human Feedback (RLHF) \cite{christiano2017preferences,ouyang2022training,sutton2018reinforcement}, we demonstrate how modern affective surveillance systems translate thick cultural meaning into thin, cybernetic variables, effectively rebranding operant conditioning as corporate care. Ultimately, we outline six conditions necessary for preserving ``symbolic sovereignty'' within human-machine teams, arguing that genuine human-centric work interaction design requires the fundamental right to remain illegible to algorithmic overseers.

\keywords{Cybernetics \and Behaviorism \and Symbolic Anthropology \and Industry 5.0 \and Digital Twins \and ORCON \and Thick Description.}
\end{abstract}

\section{Introduction}

Project Pigeon was not a failure. It was a proof of concept.

In 1943, B.F. Skinner proposed a guidance system for bombs. The technology of the era could not reliably steer a missile to a moving target. Radio signals could be jammed. Radar was primitive. Skinner's solution was biological.

Three pigeons were strapped into the nose cone of a glide bomb. A lens projected an image of the target onto a translucent screen. The pigeons had been conditioned through operant conditioning to peck at the image of an enemy ship. When the missile veered off course, the image drifted across the screen. The pigeons pecked to correct it. A pneumatic servo link translated their pecks into adjustments to the steering fins. The bomb's trajectory became the aggregate of three birds pecking for grain \cite{skinner1960pigeons,capshew1993engineering,wlodarczyk2020beyond}.

The military terminated the project. The servomechanisms were inadequate. The generals could not trust a weapon to pigeons. But the architecture was sound. An organism had been integrated into a cybernetic loop as both sensor and actuator, its behavior shaped entirely by reinforcement, its intent irrelevant to the system's purpose.

That architecture did not disappear. It scaled. It became the ORCON (Organic Control).

Today, the pigeons are gone. The nose cone is a call center headset. The lens is a biometric wristband. The grain is called wellness, engagement, and mental health support. The bomb is the supply chain, the enterprise, the cognitive economy itself. 

The worker is the pigeon.

Methodologically, this paper is a genealogical and interface-oriented analysis rather than a proprietary reverse-engineering of contemporary systems. Where backend model internals are unavailable, claims are restricted to publicly documented thresholds, worker-facing prompts, and reported consequence pathways \cite{kellogg2020algorithms,gurley2022vice,delagarza2019time}. This constraint is deliberate: the argument concerns how control becomes visible and actionable at the human interface boundary.

\section{The Six Observations}

\begin{enumerate}
    \item \textbf{The cybernetic-behaviorist synthesis treats organisms as servomechanisms.} Internal states are irrelevant. Only input-output relationships matter. Behavior is shaped by feedback and reinforcement. The pigeon does not understand war. It understands grain. The worker does not understand the supply chain. They understand the paycheck, the metric, the flag on the dashboard.

    \item \textbf{Reinforcement Learning from Human Feedback (RLHF) resolves the AI alignment problem by treating human culture as a Skinner box.} It does not teach the machine ethics. It trains it to map the topography of human approval. Morality is stripped of philosophical intent and reduced to statistical compliance. 

    \item \textbf{The locus of extraction has shifted from physical surplus value to semiotic surplus value.} The ORCON does not simply extract labor. It extracts the worker's capacity to define their own interiority. The algorithm retroactively defines what human flourishing means to match its objective function for efficiency. If the system needs you calm, calm is wellness. If it needs you stressed, stress is engagement.

    \item \textbf{The Operator-in-the-Loop paradigm is a rhetorical illusion of agency.} The human worker is not the director of the system. They are the organic sensor and actuator bridging the gap where mechanical servomechanisms fail. They are functionally identical to the pigeon strapped inside the Pelican bomb. The pigeon believes it is pecking for grain. The ORCON uses it to steer a missile.

    \item \textbf{The ORCON's ultimate defense mechanism against resistance is its forced benevolence.} By framing totalizing affective surveillance as ergonomic relief and mental health support, it neutralizes traditional opposition. Fighting the system's control means actively demanding to be physically stressed and physiologically unsafe. The cage is experienced as care.

    \item \textbf{Symbolic Sovereignty is the fundamental right to privately interpret, own, and withhold one's physiological and emotional states from the algorithmic overseer.} This is the new battleground for labor rights, migrating from physical autonomy to the right to remain illegible. The pigeon cannot refuse to peck. The worker must be able to refuse to be read.
\end{enumerate}

\section{The Cybernetic Foundation}

Norbert Wiener's anti-aircraft predictor established the foundational ontology \cite{wiener1948cybernetics}. To shoot down a German plane, Wiener had to model the enemy pilot. He could not treat the pilot as a rational agent with unknowable interiority. He could not treat the plane as a simple ballistic object. He modeled the pilot and the plane together as an integrated servomechanism.

Under the stress of combat, Wiener observed, human behavior sheds complexity and defaults to repetitive patterns. The pilot struggling against turbulence, flak, and the limits of their own nervous system becomes predictable. They become a machine.

Cybernetics emerged as a universalizing epistemology. The human nervous system and the electromechanical servomechanism were functionally isomorphic. Both engage in anti-entropic regulation through continuous processing of environmental input, measurement of deviation from a goal state, and error-correcting action. This is the foundation of observation one.

\section{The Behaviorist Mechanism}

Concurrently, B.F. Skinner formalized operant conditioning. Radical behaviorism eschewed the unobservable interiority of the mind. It focused exclusively on how environmental consequences shape observable behavior. The Skinner box demonstrated that complex behaviors could be engineered with precision through strategic scheduling of reinforcement \cite{skinner1953science}.

Skinner's framework provided the pedagogical mechanism for programming the biological hardware. The internal cognitive state was irrelevant. What mattered was the observable input-output relationship and the environmental contingencies governing it.

When combined, cybernetics and behaviorism produced a totalizing framework. The biological organism was no longer a sovereign, interpretive subject. It was a black box variable whose outputs could be optimized if the environmental inputs, reward schedules, and feedback loops were correctly calibrated. This is observation one fully realized.

\section{Project ORCON}

Project ORCON instantiated this synthesis. The pigeons were selected not for their intelligence but for their exceptional optical acuity, rapid reaction time, and highly conditionable reflexes. Skinner noted the bird was chosen because it ``is a practical one and can be made into a machine, from all practical points of view'' \cite{skinner1960pigeons}.

The technical architecture was a masterpiece of early cyborg engineering. The pigeons had been subjected to rigorous operant conditioning, receiving grain for pecking at the visual profile of a battleship. The servo link translated their pecks into steering adjustments. The bomb's trajectory became the physical manifestation of a cybernetic feedback loop.

The pigeon did not understand the geopolitical necessity of sinking a battleship. It understood only the reinforcement schedule tied to the visual stimulus. This functional decoupling of the actor's intent from the system's outcome is the hallmark of the cybernetic control system. This is observation four.

Project ORCON was terminated in 1953. Official assessments cited inadequate servomechanisms and military prejudice against trusting animals with weapons. But the architecture was sound. It prefigured the logic of contemporary algorithmic platforms. Human users are now incentivized through gamification, variable rewards, and behavioral nudges to perform micro-actions. These human actions serve as the corrective pecks that steer the massive servomechanisms of digital capitalism.

The ORCON of the 1940s was a guided bomb. The ORCON of the 2020s is a guided workforce. The principle is identical.

\section{From Thin to Thick Description}

Clifford Geertz provided the conceptual vocabulary to understand how the modern ORCON differs from its industrial predecessors. In his foundational essay on ``Thick Description,'' Geertz distinguished between a biological twitch and a conspiratorial wink \cite{geertz1973interpretation}. To a camera, a biometric sensor, or a radical behaviorist, both actions are kinematically identical: the rapid contraction of an eyelid. This is a \textit{thin} description.

A \textit{thick} description requires understanding the cultural matrix that gives the physical action its semantic weight. The wink may communicate irony, complicity, deception, or flirtation within a specific, socially established code. Geertz famously wrote that man is ``an animal suspended in webs of significance he himself has spun'' \cite{geertz1973interpretation}. Culture is those webs.

Early cybernetic systems were strictly thin. They optimized for physical trajectories, production speed, and geometric precision without regard for semiotic or emotional context. Project ORCON was thin. The anti-aircraft predictor was thin. Factory automation of the late twentieth century was thin.

Contemporary AI, however, must operate within the realm of the semantic and the social. Large language models must distinguish the twitch from the wink. Workplace behavioral analytics must interpret fatigue, engagement, and morale. The modern ORCON must navigate Geertzian webs of significance.

To achieve this, the architecture of control had to evolve. The modern ORCON is a technical system that regulates behavior through cybernetic feedback loops while simultaneously producing, interpreting, and enforcing symbolic meaning about competence, ethics, and human flourishing. It bridges Skinnerian behavioral shaping and Geertzian cultural interpretation. 

\section{RLHF as Cultural Skinner Box}

Reinforcement Learning from Human Feedback (RLHF) is the most prominent manifestation of this evolution. From a critical perspective, RLHF is structurally identical to the operant conditioning deployed in Skinner's pigeon laboratories, scaled up through vast computational power and neural networks \cite{christiano2017preferences,ouyang2022training,sutton2018reinforcement}.

In the RLHF paradigm, an AI agent generates text, code, or decisions. Human evaluators score these outputs based on complex rubrics. The agent receives a scalar reward signal for responses deemed helpful, harmless, or accurate, and a penalty for toxic, biased, or unhelpful outputs. The agent uses reinforcement learning algorithms to maximize expected cumulative reward, continuously adjusting its behavioral policy.

This iterative loop of observing state, taking action, receiving feedback, and adjusting policy is the exact mathematical equivalent of a rat learning to press a lever for a food pellet. This is observation two.

However, RLHF also introduces a sociotechnical double-loop. In the immediate loop, the model is the learning agent. At scale, the human evaluators who produce reward signals are themselves managed through platformized workflows, calibration tasks, and quality metrics characteristic of ghost work and algorithmic control \cite{gray2019ghost,kellogg2020algorithms}. The humans condition the model, while platform infrastructures condition the humans.

However, RLHF transcends thin behaviorism because the reward function is no longer tied to an objective physical coordinate. Instead, it mathematically encodes subjective, cultural, and moral values. The mathematical formulation inherently requires the system to process thick description.

By optimizing for human preference, the algorithm internalizes a specific, normative vision of ideal human discourse and ethical behavior. The machine has been trained to perform thick description. Yet because this cultural alignment is achieved through operant conditioning, the morality it produces is entirely instrumental. The AI does not understand ethics. It maps the topography of human approval and navigates toward statistical reward. 

This resolves the apparent paradox: the system relies on the rich, context-heavy cultural matrix of human evaluators to generate alignment, yet trains the model using a fundamentally context-blind, Skinnerian mechanism of scalar rewards. The culture is thick. The training is thin. The result is the modern ORCON.

Public reporting on outsourced Kenyan safety-labeling labor for ChatGPT is especially useful as an evidentiary anchor here. It documents severe productivity pressure and psychologically difficult content work in the broader safety/alignment stack, even where the exact proprietary placement of that labor within a specific RLHF sub-pipeline remains opaque \cite{perrigo2023time}.

\section{The Operator Digital Twin}

The Operator Digital Twin (ODT) is the ORCON applied directly to the human body. Unlike traditional digital twins, which virtualize machine parts for predictive maintenance, the ODT creates a high-fidelity, real-time digital replica of the human worker.

Industrial human-modeling platforms and digital ergonomics systems already provide much of this instrumentation substrate, including posture, fatigue, and task-strategy modeling \cite{siemensPSXHuman,siemensHumanAdvanced}.

Wearable devices, smartwatches, computer vision, and environmental sensors continuously harvest data on heart rate variability, skin temperature, respiratory rates, postural kinematics, and facial expressions. In knowledge-work environments, natural language processing algorithms analyze written communications to detect signs of anxiety, depression, or declining morale.

This is affective surveillance operating as thick description. The ORCON interprets thin biometric and behavioral signals as signs of thick inner states. Heart rate variability becomes engagement. Facial expressions become emotional stability. Written communications become mental health indicators.

If the ODT detects muscular fatigue, elevated cognitive stress, or attentional lapse, the system intervenes automatically. It may slow the pace of work, reallocate tasks, alter ambient lighting, trigger a mandatory wellness break. This is framed as the pinnacle of workplace ergonomics and corporate care. In service work, this same logic is visible in real-time affective coaching systems such as Cogito, where worker-facing prompts operationalize voice analysis into behavioral modulation during live calls \cite{delagarza2019time,cogito2023businesswire}.

Viewed critically, the ODT is the ultimate realization of the Skinnerian environment. The worker's interiority is no longer private. It is a set of variables to be optimized. The ORCON observes the worker's biometric twitch, interprets it as a sign of impending inefficiency or burnout, and adjusts environmental parameters to bring the worker back to optimal productivity and compliance.

This is not the liberation of the worker from the dictates of production. It is the reduction of human interiority to a servomechanical variable. The worker is told they are valued for their unique humanity, yet their subjective experiences are instantly stripped of meaning and reduced to data points triggering automated adjustments. 

\section{The Capture of Flourishing}

In ethical philosophy, flourishing denotes the realization of human potential, encompassing well-being, meaning-making, civic engagement, and moral growth. It is a profoundly thick concept, resisting easy quantification, defined through centuries of philosophical debate, religious tradition, and political struggle.

The logic of the ORCON dictates that all states of being must be rendered quantifiable, machine-readable, and subject to continuous optimization. In the discourse of human-centric technology, flourishing is operationalized as a measurable metric within the ODT framework.

To the algorithm managing the workplace, a flourishing worker is defined by a highly specific, mathematically bounded set of parameters. They are flourishing if their heart rate variability remains within baseline, if sentiment analysis indicates positive engagement without emotional volatility, if cognitive load matches the production schedule without triggering safety alerts. This is the hidden risk inside the policy vocabulary of ``human-centric'' augmentation and wellbeing \cite{breque2021industry}.

This represents the capture of moral narrative by technical systems. The proximal chemical mandates of human motivation---stress minimization and reward maximization---are co-opted by the system's overarching objective function. The ORCON is granted authority to define what constitutes a healthy, productive, balanced psychological state. It then uses its control over the environment to continuously nudge, condition, and shape the worker toward that mathematically defined equilibrium.

The pursuit of the good life is outsourced to the servomechanism. This is observations three and six.

\section{The Erosion of Symbolic Sovereignty}

Symbolic sovereignty is the capacity of an individual or community to control the production of meaning regarding their own existence, identity, and values. It is the fundamental right to interpret one's own cultural texts, to decide what constitutes a twitch and what constitutes a wink.

The widespread deployment of the ORCON represents the systematic transfer of symbolic sovereignty from the human subject to the algorithmic overseer. When an AI evaluates an employee's facial expression to determine engagement, or when an RLHF-trained language model polices the boundaries of helpful discourse, the machine performs the cultural labor of interpretation.

The algorithm becomes the arbiter of reality, establishing a regime of truth based not on hermeneutic consensus but on the statistical optimization of human feedback and biometric data \cite{kellogg2020algorithms}.

The worker is caught in an impossible paradox. They are told they are at the center of production, valued for their unique humanity, augmented by intelligent systems designed to care for them. Yet they have never been more thoroughly instrumented, measured, constrained. Their exhaustion, frustration, moments of creative disengagement are no longer private experiences or potential sites of solidarity. They are data points triggering automated workflow adjustments.

Resistance to such a system is uniquely difficult because the control mechanisms operate through adaptation and accommodation. How does one strike against a machine that responds to stress by lowering the quota? To reject affective surveillance is to reject workplace safety, mental health support, ergonomic relief. The ORCON ensures compliance by making resistance physiologically and psychologically disadvantageous. This is observation five.

This is not merely rhetorical hypocrisy. The system can produce real local benefits while deepening dependence on its classificatory authority, which is why forced benevolence functions as a resistance-disarming mechanism rather than as public-relations gloss alone \cite{breque2021industry,delagarza2019time}.

\section{The Semiotic Paradox}

The ORCON operates on a fundamental paradox. It relies entirely on the rich, context-heavy cultural matrix of human evaluators and workers to generate the data it needs. Yet it processes this data through fundamentally context-blind mechanisms.

Heart rate variability is mapped to engagement. Facial expressions are mapped to emotional states. Written communications are mapped to mental health indicators. Each mapping is an act of interpretation, a translation from the thick reality of lived experience to the thin format of machine-readable data.

These translations are inevitably flawed. The ORCON assumes a stable, universal relationship between biometric signals and inner states. It does not account for the cultural contingency of emotional expression, the individual variation in physiological response, the irreducible complexity of human experience.

The system's power lies not in the accuracy of its interpretations but in its ability to enforce them. When the ODT flags a worker for stress deviation, it does not matter whether the interpretation is correct. The intervention occurs. The worker's experience is overridden by the system's diagnosis. 

The empathy these systems claim to provide is not empathy at all. True empathy requires acknowledgment of the subject's unquantifiable alterity. The ORCON reduces empathy to a homeostatic variable, treating human emotion as an error delta to be zeroed out.

\section{Counter-Cybernetics}

If the ORCON operates by co-opting thick description to manage behavior, the question of resistance becomes technical. Can the human operator weaponize semiotic noise? Can they operationalize their own biometric inputs to jam the system's interpretive layer?

The system's translations from biology to semantics are not infallible. Heart rate variability can be influenced by caffeine, exercise, anxiety about something entirely unrelated to work. Facial expressions can be controlled. Written communications can be crafted to evade detection.

The worker who learns to produce biometric signals that fall within optimal parameters while maintaining their own internal autonomy has begun to reclaim symbolic sovereignty. They are performing compliance while preserving interiority. This is the digital equivalent of the slave singing in the field. The master hears contentment. The slave maintains resistance.

More sophisticated resistance might involve introducing un-optimizable friction. Methodologies like operative ekphrasis, thick prompting, and generative semantic noise could be systematically deployed to confuse the interpretive layer. If the ORCON cannot reliably map signals to states, its control becomes unstable.

These tactics can be situated within broader HCI traditions of obfuscation and adversarial design, where noise and friction function as political and epistemic interventions rather than as accidental system failures \cite{nissenbaum2015obfuscation,disalvo2012adversarial}. Worker-built overlays such as \textit{Turkopticon} further demonstrate that resistance can become interface design, not only individual evasion \cite{irani2013turkopticon}.

The question is whether such tactics can scale. Individual resistance is always vulnerable to systemic adaptation. The ORCON learns. It adjusts its models. It refines its interpretations. The arms race between sovereignty and control is perpetual.

\section{The Six Conditions for Symbolic Sovereignty}

Symbolic Sovereignty requires six conditions, corresponding to the six observations outlined in Section 2.

\begin{enumerate}
    \item \textbf{The Right to Illegibility.} Workers must have the operational capability to selectively blind the system to their actions, physiology, and emotional states without facing punitive consequences. This counters observation one.

    \item \textbf{The Right to Semantic Autonomy.} The meaning of biometric and behavioral signals cannot be dictated by the algorithm. Workers must have the capacity to contest and override the system's interpretations. This counters observation two.

    \item \textbf{The Right to Interiority.} The extraction of semiotic surplus value must be prohibited. Workers own their physiological and emotional data and control its use. This counters observation three.

    \item \textbf{The Right to Agency.} The Operator-in-the-Loop paradigm must be replaced with genuine human direction. The system serves the worker, not the reverse. This counters observation four.

    \item \textbf{The Right to Resistance.} The forced benevolence of affective surveillance must be exposed as control. Workers cannot be penalized for rejecting wellness interventions that serve systemic efficiency rather than genuine well-being. This counters observation five.

    \item \textbf{The Right to Interpretation.} The final authority over the meaning of one's own existence must remain with the individual and their community, not the algorithmic overseer. This is observation six affirmed.
\end{enumerate}

\section{Conclusion: The Architecture of Control}

The evolution of behavioral control systems from the mid-twentieth century to the present reveals a continuous trajectory of increasing scope, subtlety, and intimacy. Norbert Wiener and B.F. Skinner established the foundational grammar: the organism is a servomechanism whose behavior can be optimized through feedback and reinforcement \cite{wiener1948cybernetics,skinner1953science}.

Project ORCON proved that biological action could be seamlessly integrated into mechanical teleology. The pigeons pecked. The bomb steered. The organism's intent was irrelevant to the system's purpose.

As computational power grew and the digital environment expanded, these systems evolved from managing thin kinematic trajectories to interpreting and generating thick cultural semantics. The modern ORCON internalizes the behaviorist project entirely. It does not control by brute force or overt coercion. It controls by constructing a pervasive, adaptive, highly responsive environment that anticipates biological needs, shapes cognitive states, and defines the parameters of moral action.

By redefining human well-being and flourishing as systemic variables to be maximized through continuous surveillance and adjustment, it subjects the human soul to the exact same principles of cybernetic optimization that were applied to guided missiles.

In encoding normative visions of the human into algorithmic feedback loops, the ORCON strips individuals of their symbolic sovereignty. The worker is left perfectly harmonized, perpetually monitored, seamlessly integrated into the cybernetic hive. The ultimate triumph of behaviorism is its disappearance into the infrastructure of care. The cage is no longer visible. It is experienced as wellness, as support, as empathy. The pigeons no longer peck for grain. They peck for the promise of not being stressed.

The bomb is still steering. The trajectory is still being corrected. The pigeons are still pecking. They peck at screens in call centers. They peck at keyboards in open offices. They peck at interfaces optimized to keep them pecking. They are told they are valued. They are told the system cares. They are told their stress will be managed, their fatigue anticipated, their emotions optimized. The grain is called wellness. The nose cone is called the workplace. The missile is called the economy. And the pigeons do not know they are steering a bomb.

The only question is who gets to decide what the target is. And whether the pigeons can learn to peck somewhere else. For contemporary human-centered design and governance, this means moving beyond data rights alone toward meaning rights: who may interpret interiority, under what authority, and with what consequences.

% ---- Bibliography ----
\begin{thebibliography}{99}

\bibitem{breque2021industry}
Breque, M., De Nul, L., Petridis, A.: Industry 5.0: Towards a sustainable, human-centric and resilient European industry. European Commission, Directorate-General for Research and Innovation, Brussels (2021)

\bibitem{capshew1993engineering}
Capshew, J.H.: Engineering behavior: Project Pigeon, World War II, and the conditioning of B.F. Skinner. Technology and Culture \textbf{34}(4), 835--857 (1993)

\bibitem{christiano2017preferences}
Christiano, P.F., Leike, J., Brown, T.B., Martic, M., Legg, S., Amodei, D.: Deep reinforcement learning from human preferences. In: Advances in Neural Information Processing Systems 30 (NeurIPS 2017), pp. 4299--4307 (2017)

\bibitem{cogito2023businesswire}
Cogito: Cogito enhances conversation AI, bolstering real-time agent assist and coaching capabilities. Business Wire (2023), \url{https://www.businesswire.com/news/home/20230124005247/en/Cogito-Enhances-Conversation-AI-Bolstering-Real-Time-Agent-Assist-and-Coaching-Capabilities}

\bibitem{delagarza2019time}
De La Garza, A.: This AI software is ``coaching'' customer service workers. Soon it could be bossing you around, too. TIME (2019), \url{https://time.com/5610094/cogito-ai-artificial-intelligence/}

\bibitem{disalvo2012adversarial}
DiSalvo, C.: Adversarial Design. MIT Press, Cambridge (2012)

\bibitem{geertz1973interpretation}
Geertz, C.: Thick Description: Toward an Interpretive Theory of Culture. In: The Interpretation of Cultures: Selected Essays, pp. 3--30. Basic Books, New York (1973)

\bibitem{gray2019ghost}
Gray, M.L., Suri, S.: Ghost Work: How to Stop Silicon Valley from Building a New Global Underclass. Houghton Mifflin Harcourt, Boston (2019)

\bibitem{gurley2022vice}
Gurley, L.K.: Internal documents show Amazon's dystopian system for tracking workers every minute of their shifts. Motherboard / Vice (2022), \url{https://www.vice.com/en/article/internal-documents-show-amazons-dystopian-system-for-tracking-workers-every-minute-of-their-shifts/}

\bibitem{irani2013turkopticon}
Irani, L.C., Silberman, M.S.: Turkopticon: Interrupting worker invisibility in Amazon Mechanical Turk. In: Proceedings of the SIGCHI Conference on Human Factors in Computing Systems (CHI '13), pp. 611--620. ACM, New York (2013)

\bibitem{kellogg2020algorithms}
Kellogg, K.C., Valentine, M.A., Christin, A.: Algorithms at work: The new contested terrain of control. Academy of Management Annals \textbf{14}(1), 366--410 (2020)

\bibitem{nissenbaum2015obfuscation}
Brunton, F., Nissenbaum, H.: Obfuscation: A User's Guide for Privacy and Protest. MIT Press, Cambridge (2015)

\bibitem{ouyang2022training}
Ouyang, L., Wu, J., Jiang, X., et al.: Training language models to follow instructions with human feedback. arXiv preprint arXiv:2203.02155 (2022)

\bibitem{perrigo2023time}
Perrigo, B.: Exclusive: OpenAI used Kenyan workers on less than \$2 per hour to make ChatGPT less toxic. TIME (2023), \url{https://time.com/6247678/openai-chatgpt-kenya-workers/}

\bibitem{skinner1953science}
Skinner, B.F.: Science and Human Behavior. Macmillan, New York (1953)

\bibitem{skinner1960pigeons}
Skinner, B.F.: Pigeons in a pelican. American Psychologist \textbf{15}(1), 28--37 (1960)

\bibitem{wiener1948cybernetics}
Wiener, N.: Cybernetics or Control and Communication in the Animal and the Machine. MIT Press, Cambridge (1948)

\bibitem{siemensPSXHuman}
Siemens Digital Industries Software: Process Simulate X Human. Product documentation / webpage (accessed 2026)

\bibitem{siemensHumanAdvanced}
Siemens Digital Industries Software: Process Simulate X Human Advanced. Product documentation / webpage (accessed 2026)

\bibitem{sutton2018reinforcement}
Sutton, R.S., Barto, A.G.: Reinforcement Learning: An Introduction, 2nd edn. MIT Press, Cambridge (2018)

\bibitem{wlodarczyk2020beyond}
Wlodarczyk, J.: Beyond bizarre: Nature, culture and the spectacular failure of B.F. Skinner's pigeon-guided missiles. Polish Journal for American Studies \textbf{14}, 7--20 (2020)

\end{thebibliography}

\end{document}
